\appendix
\chapter*{Přílohy}
\fancyfoot[R]{\thepage}
\thispagestyle{fancy}

\addcontentsline{toc}{chapter}{Přílohy}

{\color{infoBlue}
Podle úvahy autora šířeji a hlouběji vysvětlují a dokreslují metody a výzkumné techniky uváděné v hlavním textu.

Mezi přílohy patří:
\begin{itemize}
    \item \textbf{Doplňkový obrazový materiál} – grafy, diagramy, nákresy, schémata, faksimile (opisy), mapy, plány, ukázky textů.
    \item \textbf{Některé tabulky} – dotýkají se hlavního tématu jen volně, nebo jsou to tabulky složitější a většího rozsahu.
    \item \textbf{Formuláře} použitých dotazníků, osnovy rozhovorů, pozorovací archy.
    \item \textbf{Bibliografie} zachycující literaturu příbuznou k předmětu práce, která však nebyla využita.
    \item \textbf{Popis počítačových programů}, nebo jiné výzkumné techniky.
\end{itemize}

Každá příloha začíná na nové stránce.
}

\pagebreak

\chapter*{Příloha A: Ukázka kódu}
\addcontentsline{toc}{chapter}{Příloha A: Ukázka kódu}

\lstset{
  basicstyle=\ttfamily\small,
  numbers=left,
  numberstyle=\tiny\color{gray},
  stepnumber=1,
  frame=single,
  breaklines=true,
  aboveskip=0mm,
  belowskip=0mm 
}

\begin{singlespace}
\lstinputlisting[
  firstline=1,
  lastline=30,
  caption={GitHub Actions workflow},
  label={code:example}
]{.github/workflows/latex.yaml}
\end{singlespace}
